% concatenated from english_version_trp.tex
%% 
%% Copyright 2007-2025 Elsevier Ltd
%% 
%% This file is part of the 'Elsarticle Bundle'.
%% ---------------------------------------------
%% 
%% It may be distributed under the conditions of the LaTeX Project Public
%% License, either version 1.3 of this license or (at your option) any
%% later version.  The latest version of this license is in
%%    http://www.latex-project.org/lppl.txt
%% and version 1.3 or later is part of all distributions of LaTeX
%% version 1999/12/01 or later.
%% 
%% The list of all files belonging to the 'Elsarticle Bundle' is
%% given in the file `manifest.txt'.
%% 
%% Template article for Elsevier's document class `elsarticle'
%% with harvard style bibliographic references

%\documentclass[preprint,12pt]{elsarticle}

%% Use the option review to obtain double line spacing
%% \documentclass[preprint,review,12pt]{elsarticle}

%% Use the options 1p,twocolumn; 3p; 3p,twocolumn; 5p; or 5p,twocolumn
%% for a journal layout:
\documentclass[final,1p,times]{elsarticle}
%% \documentclass[final,1p,times,twocolumn]{elsarticle}
% \documentclass[final,3p,times]{elsarticle}
%%\documentclass[final,3p,times,twocolumn]{elsarticle}
%% \documentclass[final,5p,times]{elsarticle}
%% \documentclass[final,5p,times,twocolumn]{elsarticle}

%% For including figures, graphicx.sty has been loaded in
%% elsarticle.cls. If you prefer to use the old commands
%% please give \usepackage{epsfig}

%% The amssymb package provides various useful mathematical symbols
\usepackage{amssymb}
%% The amsmath package provides various useful equation environments.
\usepackage{amsmath}
%% The amsthm package provides extended theorem environments
\usepackage{amsthm}
\usepackage{url}
\usepackage{float}
\usepackage[T1]{fontenc}
\usepackage[utf8]{inputenc}
\usepackage{enumitem}

%% The lineno packages adds line numbers. Start line numbering with
%% \begin{linenumbers}, end it with \end{linenumbers}. Or switch it on
%% for the whole article with \linenumbers.
%% \usepackage{lineno}




\newcommand{\Hcal}{\mathcal{H}}
\newcommand{\Acal}{\mathcal{A}}

\newcommand{\kstar}{\kappa^{*}\!}
\newcommand{\ddt}{\dfrac{\mathrm{d}}{\mathrm{d}\,t}}
\newcommand{\dx}{\,\mathrm{d}x}

\newtheorem{innerhyp}{(H)}
\newenvironment{hypothesis}[1]
{\renewcommand\theinnerhyp{#1}\innerhyp}
{\endinnerhyp}

\journal{ }

\numberwithin{equation}{section}
\newtheorem{theorem}{Theorem}[section]
\newtheorem{lemma}[theorem]{Lemma}
\newtheorem{corollary}[theorem]{Corollary}
\newtheorem{definition}[theorem]{Definition}
\newtheorem{proposition}[theorem]{Proposition}
\newtheorem{remark}[theorem]{Remark}
\newtheorem{example}[theorem]{Example}
\allowdisplaybreaks


\newcommand{\R}{\mathbb{R}}
\newcommand{\N}{\mathbb{N}}
\newcommand{\QT}{Q_{\mathcal{T}}}
\newcommand{\dd}{\,\mathrm{d}}
%\newcommand{\dx}{\,\mathrm{d}x}
\newcommand{\dt}{\,\mathrm{d}t}
\newcommand{\ds}{\,\mathrm{d}s}
\newcommand{\norm}[1]{\left\|#1\right\|}
\newcommand{\normsq}[1]{\left\|#1\right\|^2}
\newcommand{\inner}[2]{\left\langle#1,\,#2\right\rangle}
\newcommand{\abs}[1]{\left|#1\right|}
\newcommand{\ppos}[1]{\left[#1\right]^{+}}
\newcommand{\pneg}[1]{\left[#1\right]^{-}}
\newcommand{\ptn}{\partial_t}
\newcommand{\Kd}{K_d(\theta)}
\newcommand{\Kr}{K_r(\theta)}

\begin{document}
	
	\begin{frontmatter}
		
		%% Title, authors and addresses
		
		%% use the tnoteref command within \title for footnotes;
		%% use the tnotetext command for theassociated footnote;
		%% use the fnref command within \author or \affiliation for footnotes;
		%% use the fntext command for theassociated footnote;
		%% use the corref command within \author for corresponding author footnotes;
		%% use the cortext command for theassociated footnote;
		%% use the ead command for the email address,
		%% and the form \ead[url] for the home page:
		%% \title{Title\tnoteref{label1}}
		%% \tnotetext[label1]{}
		%% \author{Name\corref{cor1}\fnref{label2}}
		%% \ead{email address}
		%% \ead[url]{home page}
		%% \fntext[label2]{}
		%% \cortext[cor1]{}
		%% \affiliation{organization={},
			%%             addressline={},
			%%             city={},
			%%             postcode={},
			%%             state={},
			%%             country={}}
		%% \fntext[label3]{}
		
		\title{Maximum principle and local stability for a class of coupled nonlinear thermo--reaction--phase systems} %% Article title
		
		%% use optional labels to link authors explicitly to addresses:
		%% \author[label1,label2]{}
		%% \affiliation[label1]{organization={},
			%%             addressline={},
			%%             city={},
			%%             postcode={},
			%%             state={},
			%%             country={}}
		%%
		%% \affiliation[label2]{organization={},
			%%             addressline={},
			%%             city={},
			%%             postcode={},
			%%             state={},
			%%             country={}}
		
		
		
		%% Author affiliation
		\author[1]{Gossrin Jean-Marc Bomisso} \ead{bogojm@yahoo.fr}
		\author[1]{Ali Ouattara Kouma} \ead{benzeus96@gmail.com} 
		\author[1]{Marie Esther Anass\'e} \ead{hadassamarie43@gmail.com}
		
		
		\address[1]{Universit\'e Nangui Abrogoua, UFR Sciences Fondamentales et Appliqu\'ees, 02 BP 801 Abidjan 02, C\^ote d'Ivoire}
		
		
		%% Abstract
		\begin{abstract}
			We study a nonlinear coupled system of partial differential equations arising from thermo--reaction--phase models. The system combines a heat diffusion equation, temperature-dependent chemical reactions of Arrhenius type, and a phase variable, and is formulated as a strongly coupled parabolic problem with homogeneous Neumann boundary conditions.
			We first establish a maximum principle ensuring the positivity of the temperature on a suitable time interval, as well as the invariance of the physically admissible domain. In particular, we prove that the internal variables remain in the interval $[0,1]$.
			We then analyse the asymptotic behaviour of the system in the free regime, that is, in the absence of external forcing. By introducing a relative energy functional and exploiting the structure of the coupling terms, we obtain local asymptotic stability of a homogeneous stationary state.
			The model belongs to a broader class of coupled diffusion--reaction--phase systems.
		\end{abstract}
		
		
		%% Keywords
		\begin{keyword}
			Thermo-reaction system \sep maximum principle \sep stability analysis \sep Arrhenius kinetics
			
			\MSC[2020] 35K57 \sep 35B50 \sep 35B35 \sep 80A32
			
			
			
		\end{keyword}
		
	\end{frontmatter}
	
	\section{Introduction}
	
	Heat-sensitive materials constitute a class of engineering materials whose physical and chemical properties evolve significantly, and sometimes irreversibly, in response to a thermal stimulus. Among the most emblematic phenomena in this class, thermochromism refers to the modification of the optical absorption spectrum of a material under the effect of heat, leading to a macroscopically observable colour change. This phenomenon results from a fine coupling between heat transport, molecular chemical transformations and phase reorganisation, making it both a rich object of study from the modelling perspective and a delicate one from the mathematical analysis point of view. From an applicative standpoint, such materials appear in sectors as varied as intelligent food packaging, technical textiles, safety paints, passive display devices, and thermal protection of electronic components. A remarkable synthesis of the physical and chemical mechanisms underlying thermochromism in polymers is provided by Seeboth et al.~\cite{SeebothEtal2014}, who distinguish in particular the mechanisms of thermal isomerisation, supramolecular reorganisation, and proton transfer. MacLaren and White~\cite{MacLarenWhite2005} have established design rules for reversible thermochromic mixtures based on leuco dyes, whose decolouration-recolouration kinetics constitutes precisely the central mechanism of the model studied here. More generally, the notion of stimuli-responsive material has been highlighted by Stuart et al.~\cite{StuartEtal2010}, who survey the emerging applications of polymers capable of modifying their structure in response to external stimuli such as temperature, pH, or light. This unifying framework shows that the thermal variable should no longer be regarded as a mere external parameter, but as a dynamic quantity coupled to the internal state variables of the material.
	
	From this perspective, the mathematical description of these phenomena naturally leads to strongly coupled nonlinear systems of partial differential equations. Let $\Omega \subset \R^N$ ($N \geq 1$) be a bounded domain with boundary $\partial\Omega$ of class $C^{1,1}$, and let $\mathcal{T} > 0$ be a finite time horizon. We denote $\QT = \Omega \times (0,\mathcal{T})$. The system studied in this work couples three unknowns, namely the absolute temperature $\theta$ (in kelvin), the active chemical fraction $c$ representing the proportion of the coloured form of the material, and the macroscopic phase variable $\phi$ describing the visibility state of the material. These three unknowns are governed by
	\begin{equation}
		\label{eq:sys}
		\left\{
		\begin{aligned}
			\rho c_p\,\ptn\theta - \nabla\cdot\bigl(k(\phi)\nabla\theta\bigr)
			&= H_{\mathrm{ext}}(x,t) + \alpha L_c\,\ptn c + \alpha L_\phi\,\ptn\phi
			&& \text{in } \QT, \\[4pt]
			\ptn c - d_c\,\Delta c
			&= -\beta K_d(\theta)\,c + \gamma K_r(\theta)\,(1-c)
			&& \text{in } \QT, \\[4pt]
			\tau_\phi\,\ptn\phi - \varepsilon^2\Delta\phi + F'(\phi)
			&= \lambda(c-\phi)
			&& \text{in } \QT,
		\end{aligned}
		\right.
	\end{equation}
	complemented by the homogeneous Neumann boundary conditions
	\begin{equation}
		\label{eq:Neumann}
		\frac{\partial\theta}{\partial\mathbf{n}} = 0,\quad
		\frac{\partial c}{\partial\mathbf{n}} = 0,\quad
		\frac{\partial\phi}{\partial\mathbf{n}} = 0
		\quad\text{on } \partial\Omega\times(0,\mathcal{T}),
	\end{equation}
	and by the initial conditions
	\begin{equation}
		\label{eq:CI}
		\theta(\cdot,0) = \theta_0,\quad
		c(\cdot,0) = c_0,\quad
		\phi(\cdot,0) = \phi_0
		\quad\text{in } \Omega.
	\end{equation}
	The model parameters are as follows: $\rho$ denotes the mass density, $c_p$ the specific heat capacity, $k(\phi)$ the phase-dependent thermal conductivity, $H_{\mathrm{ext}}(x,t)$ an external heat source, $L_c$ and $L_\phi$ the thermal coefficients associated respectively with the chemical transformation and the phase transition, $\alpha$ the internal thermal coupling intensity parameter, $d_c$ the chemical diffusion coefficient, $\beta$ and $\gamma$ the respective intensities of the decolouration and recolouration mechanisms, $\tau_\phi$ the phase relaxation time, $\varepsilon$ the interface parameter, and $\lambda$ the relaxation coefficient coupling $c$ to $\phi$.
	
	The kinetic rates are assumed to be of Arrhenius type,
	\begin{equation}
		\label{eq:Arrhenius}
		K_d(\theta) = A_d \exp\!\left(-\frac{E_d}{R\theta}\right),
		\qquad
		K_r(\theta) = A_r \exp\!\left(-\frac{E_r}{R\theta}\right),
		\quad \theta > 0,
	\end{equation}
	where $A_d, A_r > 0$ are pre-exponential factors, $E_d, E_r > 0$ are activation energies, and $R > 0$ is the universal gas constant. For the phase variable, we choose the double-well potential
	\begin{equation}
		\label{eq:F}
		F(\phi) = \tfrac{1}{4}\phi^2(1-\phi)^2, \qquad
		F'(\phi) = \tfrac{1}{2}\phi(1-\phi)(1-2\phi),
	\end{equation}
	whose two minima at $\phi = 0$ and $\phi = 1$ correspond to the two stable states of the material.
	
	In order to place the analysis within a rigorous framework, we assume that the following hypotheses are satisfied. They are grouped into two families according to their scope within the paper.
	
	\smallskip
	\noindent\textit{General hypotheses (valid throughout the paper).}
	
	\begin{itemize}
		\item[\textbf{(H1)}] \textit{(Thermal conductivity)}
		$k\in C^1(\mathbb R)$ and there exist two constants $0<k_*\le k^*<\infty$ such that
		$k_*\le k(s)\le k^*$ for all $s\in\mathbb R$.
		
		\item[\textbf{(H2)}] \textit{(Kinetic rates)}
		$K_d,K_r\in C^1(0,\infty)$, with $K_d,K_r\ge 0$. Moreover, for every fixed $\theta_*>0$, $K_d$ and $K_r$ are bounded on $[\theta_*,\infty)$. The Arrhenius expressions \eqref{eq:Arrhenius} satisfy these conditions as soon as $\theta$ is uniformly bounded below.
		
		\item[\textbf{(H3)}] \textit{(Phase potential)}
		$F\in C^2(\mathbb R)$, $F\ge 0$, is given by \eqref{eq:F}.
		
		\item[\textbf{(H4)}] \textit{(Parameters)}
		$\rho,c_p,d_c,\tau_\phi,\varepsilon,\lambda,\beta,\gamma>0$ and $\alpha, \, L_c, \, L_{\phi} \, \ge 0$.
		
		\item[\textbf{(H5)}] \textit{(Thermal source)}
		We write $S = H_{\mathrm{ext}} + \alpha L_c\,\ptn c + \alpha L_\phi\,\ptn\phi$.
		We assume
		\[
		H_{\mathrm{ext}} \in L^2(0,\mathcal{T};\,L^2(\Omega)),
		\qquad
		S \in L^\infty(Q_{\mathcal{T}}),
		\qquad
		\exists\, C_0 \geq 0 \colon S(x,t) \geq -C_0
		\quad \text{a.e. in } Q_{\mathcal{T}}.
		\]
		
		\item[\textbf{(H6)}] \textit{(Initial data)}
		$\theta_0\in L^2(\Omega)$ with $\theta_0\ge\theta_*>0$ a.e. in $\Omega$ and $\theta_0 \in L^\infty(\Omega)$;
		$c_0\in H^1(\Omega)$ with $0\le c_0\le 1$ a.e. in $\Omega$;
		$\phi_0\in H^1(\Omega)$ with $0\le\phi_0\le 1$ a.e. in $\Omega$.
		
	\end{itemize}
	
	\smallskip
	\noindent\textit{Additional hypotheses (required for Section~\ref{sec:stab} only).}
	
	\begin{itemize}
		\item[\textbf{(H7)}] \textit{(Absence of external source)}
		$H_{\mathrm{ext}}=0$.
		
		\item[\textbf{(H8)}] \textit{(Homogeneous stationary state)}
		There exists $(\bar\theta,\bar c,\bar\phi)$ satisfying $\bar\theta>0$, $0\le\bar c\le 1$, $0\le\bar\phi\le 1$, and the equilibrium relations
		\[
		-\beta K_d(\bar\theta)\bar c+\gamma K_r(\bar\theta)(1-\bar c)=0,
		\qquad
		F'(\bar\phi)=\lambda(\bar c-\bar\phi).
		\]
		
		\item[\textbf{(H9)}] \textit{(Local stability of the potential)}
		$F''(\bar\phi)>0$.
		
		\item[\textbf{(H10)}] \textit{(Zero-mean condition and weak coupling)} 
		We assume that
		\[
		\int_\Omega (\theta_0 - \bar{\theta})\,dx = 0,
		\]
		and that the parameter $\alpha$ satisfies $0 \leq \alpha < \alpha_0$, where
		\begin{equation}
			\label{eq:alpha0}
			\alpha_0 = \min\left\{
			\frac{\sqrt{\rho c_p}}{\sqrt{2}\,L_c},\;
			\frac{\sqrt{\rho c_p\,(2m_F+\tau_\phi)}}{\sqrt{2}\,L_\phi},\;
			\frac{\rho c_p\,\sqrt{d_c\,k_*}}{2\,k^*\,L_c},\;
			\frac{\rho c_p\,\varepsilon\,\sqrt{k_*}}{2\,k^*\,L_\phi}
			\right\}.
		\end{equation}
		
	\end{itemize}
	
	
	In this framework, the system \eqref{eq:sys} fits into several classical streams of the literature. On the one hand, the presence of Arrhenius laws brings it close to models arising in combustion theory, where the exponential dependence on $\tfrac{1}{\theta}$ produces well-known analytical difficulties, particularly near $\theta=0$; in this regard, the monograph of Bebernes and Eberly~\cite{BebernesEberly1989} remains an important reference for the rigorous study of nonlinear parabolic problems involving such reaction terms. On the other hand, the reaction-diffusion term governing the evolution of $c$ belongs to the broader topic of global existence and mass control in coupled systems, as illustrated by the survey of Pierre~\cite{PierreM2010}. The variable $\phi$, in turn, falls within the phase-field theory initiated by Cahn and Hilliard~\cite{CahnHilliard1958} and Allen and Cahn~\cite{AllenCahn1979}, and widely developed since. In this direction, Elliott and Zheng~\cite{ElliottZheng1986} established fundamental results for the Cahn--Hilliard equation, while Evans, Soner and Souganidis~\cite{EvansEtal1992} shed light on the connection between the Allen--Cahn equation and motion by mean curvature. The monograph of Miranville~\cite{MiranvilleBook2019} provides a comprehensive overview of the analysis of Cahn--Hilliard and Allen--Cahn type systems. Added to this family of models is the central reference of the Caginalp model~\cite{Caginalp1986}, which couples temperature and order parameter via a thermal relaxation term; Colli and Sprekels~\cite{ColliSprekels1995} subsequently studied several variants of this framework by proving the existence of global weak solutions. The Penrose--Fife model~\cite{PenroseFife1990}, which is thermodynamically consistent, has deeply influenced the analysis of systems with temperature singularities, and the works of Kenmochi and Niezg\'odka~\cite{KenmochiNiezgodka1994} illuminate the difficulties related to nonlinear dependencies on $\tfrac{1}{\theta}\cdotp$ Finally, concerning the asymptotic behaviour, the general theory of dissipative dynamical systems developed by Temam~\cite{TemamBook1997} and the results of Miranville and Zelik~\cite{MiranvilleZelik2004} on exponential attractors constitute a natural reference framework.
	
	The present work stands at the confluence of these different streams, but distinguishes itself from them by a richer and more strongly coupled structure. Indeed, the system \eqref{eq:sys} introduces a chemical variable $c$ subject to a reversible Arrhenius-type kinetics, with coexistence of decolouration ($K_d$) and recolouration ($K_r$) mechanisms, and couples this variable to an Allen--Cahn type phase dynamics via the term $\lambda(c-\phi)$. Moreover, the heat equation is affected by time-derivative coupling terms, $\alpha L_c\,\ptn c$ and $\alpha L_\phi\,\ptn\phi$, which strengthens the nonlinearity of the problem and complicates the construction of energy estimates. To the best of our knowledge, this architecture does not appear in classical models of thermochromism. The first is the singularity of the Arrhenius rates near $\theta=0$, which requires guaranteeing strict positivity of the temperature even before the equations on $c$ can be written. The second is the strongly non-symmetric nature of the coupling between the three equations, since each possesses its own structure: nonlinear diffusion for heat, semilinear reaction-diffusion for $c$, and an Allen--Cahn type equation for $\phi$. The third is the hierarchical interdependence of the a priori bounds: the positivity of $\theta$ conditions the well-definedness of $K_d(\theta)$ and $K_r(\theta)$, which in turn conditions the control of $c$, which itself appears in the control of $\phi$ via the forcing term $\lambda(c-\phi)$. More recent contributions have also investigated multiphysics phase-field systems 
	involving thermal and reactive couplings. In this direction, Miranville et al. in~\cite{MiranvilleQuintanillaSaoud2020} studied the asymptotic behaviour 
	of a coupled Cahn-Hilliard/Allen-Cahn system with temperature, establishing the 
	existence of exponential and finite-dimensional global attractors under both the 
	classical Fourier heat law and a type~III heat conduction law. This framework was 
	further extended by Makki et al. (see~\cite{MakkiMiranvilleSaoud2020}), 
	who incorporated singular potentials of logarithmic type into the same class of 
	thermally coupled phase-field systems, and derived additional regularity results 
	including a strict separation property from the pure phases. From a different 
	perspective, Ntsokongo in~\cite{Ntsokongo2023} analysed an Allen-Cahn type equation 
	nonlinearly coupled with a temperature equation and motivated by applications in 
	chemistry, proving well-posedness as well as the existence of finite-dimensional 
	global and exponential attractors.
	
	The analysis carried out here rests on a two-step strategy. In the first step, we establish a maximum principle adapted to the system, which ensures the preservation of the physically admissible domain and, in particular, the positivity of the temperature on an appropriate time interval. This result is fundamental since it guarantees that the Arrhenius kinetic laws remain well-defined and bounded, while establishing, as a mathematical consequence, that $c$ and $\phi$ remain in the interval $[0,1]$ for all time. In the second step, we study the local asymptotic stability of a homogeneous stationary state by means of a suitably chosen relative energy functional. The idea consists of linearising the system around this state, testing the perturbed equations by the unknowns themselves, and then combining the resulting estimates with Poincaré and Gronwall type inequalities in order to obtain exponential decay under appropriate coercivity assumptions.  The model studied in the present work belongs 
	to this broader class of coupled diffusion--reaction--phase systems, but 
	distinguishes itself by the simultaneous presence of Arrhenius-type chemical 
	kinetics, a reversible reaction-diffusion equation for the active fraction, and 
	an Allen--Cahn phase dynamics, all three of which interact through the thermal 
	variable.
	
	The organisation of the paper is as follows. Section~\ref{sec:prelim} gathers the preliminary tools needed for the analysis, in particular the positive and negative parts of a measurable function, the Gronwall lemma, and the abstract parabolic maximum principle. Section~\ref{sec:max} is devoted to the maximum principle for the system, establishing sequentially the positivity of $\theta$, then the invariance of $[0,1]$ for $c$, and then for $\phi$. Section~\ref{sec:stab} develops the local asymptotic stability analysis. Finally, a general conclusion synthesises the main results obtained, puts their scope in perspective, and outlines directions for future research, notably with a view to a global analysis and possible numerical validations.
	% ============================================================
	\section{Preliminary tools}
	\label{sec:prelim}
	% ============================================================
	
	We recall in this section the definitions and abstract results that will be used systematically in the proofs.
	
	\begin{definition}[\textbf{Sub-solution and super-solution} \cite{EvansBook2010}]
		Let $f : \QT \times \R \to \R$ be a function that is Lipschitz continuous with respect to its last variable. We say that $\underline{u}$ (resp. $\overline{u}$) is a \emph{sub-solution} (resp. \emph{super-solution}) of the parabolic problem with Neumann conditions
		\[
		\ptn u - \nabla\cdot(a(x)\nabla u) = f(x,t,u), \quad
		\frac{\partial u}{\partial \mathbf{n}} = 0 \text{ on } \partial\Omega,\quad
		u(\cdot,0) = u_0,
		\]
		if $\underline{u}$ (resp. $\overline{u}$) possesses the same regularity properties as a weak solution and satisfies
		\[
		\ptn \underline{u} - \nabla\cdot(a\nabla \underline{u}) \leq f(x,t,\underline{u})
		\quad (\text{resp. } 	\ptn \overline{u} - \nabla\cdot(a\nabla \overline{u}) \geq f(x,t,\overline{u})) \quad \text{a.e. in } \QT,
		\]
		\[
		\frac{\partial \underline{u}}{\partial\mathbf{n}} \leq 0
		\quad \left( \text{resp. } 	\frac{\partial \overline{u}}{\partial\mathbf{n}} \geq 0 \right) \quad \text{on } \partial\Omega,
		\]
		\[
		\underline{u}(x,0) \leq u_0(x)
		\quad (\text{resp. } \overline{u}(x,0) \geq u_0(x)).
		\]
	\end{definition}
	
	\begin{definition}[\textbf{Positive and negative parts} \cite{Ziemer}]
		For any measurable function $u : \Omega \to \R$, we define
		\[
		\ppos{u}(x) = \max(u(x), 0), \qquad \pneg{u}(x) = \max(-u(x), 0).
		\]
		We have $$u = \ppos{u} - \pneg{u}, \; |u| = \ppos{u} + \pneg{u}, \; \ppos{u} \cdot \pneg{u} = 0.$$
		If $u \in H^1(\Omega)$, then $\ppos{u}, \pneg{u} \in H^1(\Omega)$ with
		\[
		\nabla\ppos{u} = \mathbf{1}_{\{u>0\}}\nabla u, \qquad
		\nabla\pneg{u} = -\mathbf{1}_{\{u<0\}}\nabla u
		\quad\text{a.e. in } \Omega.
		\]
	\end{definition}
	
	\begin{lemma}[\textbf{Gronwall, integral form} \cite{EvansBook2010}]
		\label{lem:Gronwall}
		Let $\psi : [0,\mathcal{T}] \to \R_+$ be an absolutely continuous function satisfying
		\[
		\psi(t) \leq a + b\int_0^t \psi(s)\ds
		\]
		for almost every $t \in (0,\mathcal{T})$, with $a, b \geq 0$. Then $\psi(t) \leq a\,e^{bt}$ for all $t \in [0,\mathcal{T}]$.
	\end{lemma}
	
	\begin{theorem}[\textbf{Parabolic maximum principle} \cite{EvansBook2010}]
		\label{thm:max_par}
		Let $u$ be a weak solution of
		$\ptn u - \nabla\cdot(a(x)\nabla u) = f(x,t,u)$
		in $\QT$, with $a \geq a_0 > 0$ and $f$ uniformly Lipschitz in $u$. If $\underline{u}$
		is a sub-solution with $u(\cdot,0) \geq \underline{u}(\cdot,0)$ a.e., then
		$u \geq \underline{u}$ a.e. in $\QT$. Likewise, if $\overline{u}$ is a super-solution with
		$u(\cdot,0) \leq \overline{u}(\cdot,0)$, then $u \leq \overline{u}$ a.e. in $\QT$.
	\end{theorem}
	
	
	% ============================================================
	\section{Maximum principle for the system}
	\label{sec:max}
	% ============================================================
	The analysis is carried out under the assumption that the system admits a sufficiently regular solution on the considered time interval.\\
\indent	We establish in this section the domain invariance properties for each of the three unknowns of the system~\eqref{eq:sys}-\eqref{eq:CI}, under hypotheses \textbf{(H1)}--\textbf{(H6)}. The results are obtained sequentially: the positivity of $\theta$ is established first, as it conditions the definition and sign of the kinetic rates, which are themselves necessary for the proofs concerning $c$ and $\phi$.
	
	\subsection{Positivity of the temperature}
	\label{sub:theta}
	
	\begin{theorem}[\textbf{Maximum principle for} $\theta$]
		\label{thm:theta_pos}
		Under hypotheses \textup{\textbf{(H1)}}-\textup{\textbf{(H6)}}, any solution $\theta$ of the system \eqref{eq:sys}-\eqref{eq:CI} satisfies, for almost every $(x,t) \in \QT$,
		\[
		\theta(x,t) \geq \underline{\theta}(t) > 0,
		\]
		where $\underline{\theta}(t) = \theta_* - \dfrac{C_0}{\rho c_p}\,t$, and this holds for all
		$t \in [0, T_0)$, with
		\[
		T_0 = \frac{\rho c_p\,\theta_*}{C_0} > 0
		\qquad\bigl(\text{with the convention } T_0 = +\infty \text{ if } C_0 = 0\bigr).
		\]
	\end{theorem}
	
	\begin{proof}
		We introduce the affine function
		\[
		\underline{\theta}(t) = \theta_* - \frac{C_0}{\rho c_p}\,t,
		\]
		which is strictly positive on $[0,T_0)$. This function (independent of $x$) is a sub-solution of the first equation of the system. Indeed, since $\nabla\underline{\theta} = 0$, the diffusion term vanishes. The time derivative is $\ptn\underline{\theta} = -\frac{C_0}{\rho c_p}$, so that:
		\[
		\rho c_p\,\ptn\underline{\theta} - \nabla\cdot(k(\phi)\nabla\underline{\theta})
		= -C_0 \leq H_{\mathrm{ext}} + \alpha L_c\,\ptn c + \alpha L_\phi\,\ptn\phi
		\quad\text{a.e. in } \QT,
		\]
		by hypothesis \textbf{(H4)}. Moreover, $\underline{\theta}(0) = \theta_* \leq \theta_0(x)$
		a.e. in $\Omega$, and $\nabla\underline{\theta}\cdot\mathbf{n} = 0$ on $\partial\Omega$.
		
		Set $w = \theta - \underline{\theta}$. By linearity, $w$ satisfies
		\[
		\rho c_p\,\ptn w - \nabla\cdot(k(\phi)\nabla w)
		= S(x,t) + C_0 \geq 0 \quad\text{a.e. in } \QT,
		\]
		with $\nabla w\cdot\mathbf{n} = 0$ on $\partial\Omega$ and $w(\cdot,0) = \theta_0 - \theta_* \geq 0$.
		
		We multiply the equation on $w$ by $-\pneg{w} \leq 0$ and integrate over $\Omega$. Using
		the identity $\ptn w \cdot \pneg{w} = -\frac12 \ptn ([\pneg{w}]^2)$ and integration by parts,
		\[
		\int_\Omega \nabla\cdot(k(\phi)\nabla w)\cdot(-\pneg{w})\dx
		= -\int_\Omega k(\phi)|\nabla\pneg{w}|^2\dx
		\leq 0,
		\]
		we obtain:
		\[
		\frac{\rho c_p}{2}\frac{\dd}{\dt}\norm{\pneg{w}(t)}_{L^2(\Omega)}^2
		+ k_*\norm{\nabla\pneg{w}(t)}_{L^2(\Omega)}^2 \leq 0.
		\]
		Hence $t \mapsto \norm{\pneg{w}(t)}_{L^2}(\Omega)^2$ is non-increasing. Since
		$w(\cdot,0) \geq 0$ implies $\pneg{w}(\cdot,0) = 0$, we conclude:
		\[
		\norm{\pneg{w}(t)}_{L^2(\Omega)}^2 = 0 \quad\text{for all } t \in [0,T_0),
		\]
		whence $w \geq 0$, i.e., $\theta(x,t) \geq \underline{\theta}(t) > 0$ a.e. in $\QT$.
		
		In particular, for every $\mathcal T<T_0$, we set
		\begin{equation}
			\label{eq3.1}
			\underline{\theta}_{\mathcal T} =\min_{t\in[0,\mathcal T]}\underline{\theta}(t)
			=\theta_*-\frac{C_0}{\rho c_p}\,\mathcal T>0.
		\end{equation}
		Since $\underline{\theta}$ is decreasing on $[0,\mathcal T]$, we deduce that
		\[
		\theta(x,t)\ge \underline{\theta}_{\mathcal T}\quad \text{a.e. in }Q_{\mathcal T}.
		\]
	\end{proof}
	
	\begin{remark}
		\label{rem:theta_global}
		When $C_0 = 0$ (everywhere non-negative total source), the sub-solution $\underline{\theta}(t) = \theta_*$
		is constant and the result holds globally in time. This situation covers the case
		$H_{\mathrm{ext}} \geq 0$, $L_c \geq 0$, $L_\phi \geq 0$ (exothermic exchanges).
	\end{remark}
	
	\begin{corollary}[\textbf{Bounds on the Arrhenius rates}]
		\label{cor:Arrhenius}
		Under the hypotheses of Theorem~\ref{thm:theta_pos}, for every $\mathcal{T} < T_0$, there
		exists $M > 0$ such that
		\[
		0 \leq K_d(\theta(x,t)) \leq M, \quad 0 \leq K_r(\theta(x,t)) \leq M
		\quad\text{a.e. in } \QT.
		\]
		As a consequence, the nonlinearities of the equation on $c$ are uniformly bounded and the system is well-defined at the level of weak formulations on $[0,\mathcal{T}]$.
	\end{corollary}
	
	\begin{proof}
		Assume $\mathcal{T} < T_0$. By Theorem~\ref{thm:theta_pos},
		there exists $\underline{\theta}_{\mathcal{T}} > 0$ such that
		$\theta(x,t) \geq \underline{\theta}_{\mathcal{T}}$ a.e.\ in
		$Q_{\mathcal{T}}$.
		
		For the upper bound, we set under hypothesis \textbf{(H4)}:
		\[
		\overline{\theta}(t)
		= \norm{\theta_0}_{L^\infty(\Omega)}
		+ \dfrac{t}{\rho c_p}\norm{S}_{L^\infty(Q_{\mathcal{T}})}.
		\]
		Since $\overline\theta$ is independent of $x$, its diffusion term
		vanishes, and
		\[
		\rho c_p\,\partial_t\overline\theta
		- \nabla\cdot\bigl(k(\phi)\nabla\overline\theta\bigr)
		= \norm{S}_{L^\infty(Q_{\mathcal{T}})}
		\geq S(x,t)
		\quad\text{a.e. in } Q_{\mathcal{T}}.
		\]
		Moreover $\dfrac{\partial\theta}{\partial\mathbf{n}} = 0$ on
		$\partial\Omega\times(0,\mathcal{T})$, and by \textbf{(H5)},
		$\overline\theta(0) = \norm{\theta_0}_{L^\infty(\Omega)} \geq \theta_0(x)$
		a.e.\ in $\Omega$. Thus $\overline\theta$ is a super-solution. Setting $z = \overline\theta - \theta$
		and multiplying the inequality satisfied by $z$ by $-\pneg{z} \leq 0$,
		the same argument as in the proof of Theorem~\ref{thm:theta_pos} gives
		\[
		\frac{\rho c_p}{2}\frac{d}{dt}\norm{\pneg{z}(t)}_{L^2(\Omega)}^2
		+ k_*\norm{\nabla\pneg{z}(t)}_{L^2(\Omega)}^2 \leq 0.
		\]
		The monotone decrease of $t\mapsto\norm{\pneg{z}(t)}_{L^2(\Omega)}^2$ and the condition
		$z(\cdot,0)\geq 0$ imply $\pneg{z} = 0$, i.e.,
		\[
		\theta(x,t) \leq \overline\theta(t)
		\leq \norm{\theta_0}_{L^\infty(\Omega)}
		+ \frac{\mathcal{T}}{\rho c_p}\norm{S}_{L^\infty(Q_{\mathcal{T}})}
		= M'
		\quad\text{a.e. in } Q_{\mathcal{T}}.
		\]
		We therefore have
		$0 < \underline{\theta}_{\mathcal{T}} \leq \theta(x,t) \leq M'$
		a.e.\ in $Q_{\mathcal{T}}$.
		Since $[\underline{\theta}_{\mathcal{T}}, M']$ is a compact subset of
		$\,(0,\infty)$ and $K_d$ and $K_r$ are continuous on $\mathbb{R}_+$
		by \textbf{(H2)}, there exists $M > 0$ such that
		\[
		0 \leq K_d(\theta(x,t)) \leq M,
		\qquad
		0 \leq K_r(\theta(x,t)) \leq M
		\quad\text{a.e. in } Q_{\mathcal{T}}.
		\]
	\end{proof}
	
	\subsection{Invariance of $[0,1]$ for the chemical fraction $c$}
	\label{sub:c}
	
	\begin{theorem}[\textbf{Confinement of} $c$ \textbf{in} {$[0,1]$}]
		\label{thm:c_bounds}
		Under hypotheses \textup{\textbf{(H1)}}-\textup{\textbf{(H6)}} and on the interval
		$[0, T_0)$ provided by Theorem~\ref{thm:theta_pos}, any solution $c$ of the
		system~\eqref{eq:sys}-\eqref{eq:CI} satisfies
		\[
		0 \leq c(x,t) \leq 1 \quad\text{a.e. in } \QT.
		\]
	\end{theorem}
	
	\begin{proof}
		\textbf{Lower bound ($c \geq 0$).}\\
		We multiply the second equation of \eqref{eq:sys} by $-\pneg{c} \leq 0$ and integrate
		over $\Omega$. The time term gives
		$\displaystyle\int_\Omega \ptn c\cdot(-\pneg{c})\dx = \frac12\frac{\dd}{\dt}\norm{\pneg{c}}_{L^2(\Omega)}^2$,
		and the diffusion term, after integration by parts, gives
		$d_c\norm{\nabla\pneg{c}}_{L^2(\Omega)}^2 \geq 0$. For the reaction term, on the set
		$\{c < 0\}$ where $\pneg{c} = -c > 0$:
		\begin{align*}
			(-\beta K_d(\theta)\,c + \gamma K_r(\theta)(1-c))\cdot(-\pneg{c})
			&= (-\beta K_d \cdot(-\pneg{c}) + \gamma K_r(1+\pneg{c}))\cdot(-\pneg{c}) \\
			&= -\beta K_d(\pneg{c})^2 - \gamma K_r\pneg{c} - \gamma K_r(\pneg{c})^2 \leq 0,
		\end{align*}
		since $K_d, K_r \geq 0$ (Corollary~\ref{cor:Arrhenius}) and $\pneg{c} \geq 0$. We obtain:
		\[
		\frac12\frac{\dd}{\dt}\norm{\pneg{c}(t)}_{L^2(\Omega)}^2 + d_c\norm{\nabla\pneg{c}(t)}_{L^2(\Omega)}^2 \leq 0.
		\]
		The monotone decrease of $t \mapsto \norm{\pneg{c}(t)}_{L^2(\Omega)}^2$ and the initial condition
		$c_0 \geq 0$ conclude.
		
		\textbf{Upper bound ($c \leq 1$).}\\
		Set $v = c - 1$, which satisfies
		$\ptn v - d_c\Delta v = -(\beta K_d + \gamma K_r)\,v - \beta K_d$,
		with $v(\cdot,0) = c_0 - 1 \leq 0$ and homogeneous Neumann conditions. We multiply by $\ppos{v} \geq 0$.
		The reaction term on $\{v > 0\} = \{c > 1\}$ equals
		$[-(\beta K_d+\gamma K_r)(\ppos{v})^2 - \beta K_d \ppos{v}] \leq 0$, hence:
		\[
		\frac12\frac{\dd}{\dt}\norm{\ppos{v}(t)}_{L^2(\Omega)}^2 + d_c\norm{\nabla\ppos{v}(t)}_{L^2(\Omega)}^2 \leq 0.
		\]
		The initial condition $\ppos{v(\cdot,0)} = 0$ concludes.
	\end{proof}
	
	\begin{remark}
		The proof of the upper bound $c \leq 1$ uses only $K_d, K_r \geq 0$, $\beta,\gamma \geq 0$
		and $c_0 \leq 1$. It does not directly require the positivity of $\theta$, which enters
		only indirectly to ensure that $K_d(\theta)$ and $K_r(\theta)$ are well-defined. The
		lower bound $c \geq 0$, in turn, uses $\gamma K_r(\theta) \geq 0$, which is automatic
		as soon as $\theta > 0$. Both proofs are therefore logically conditioned by
		Theorem~\ref{thm:theta_pos}.
	\end{remark}
	
	\subsection{Invariance of $[0,1]$ for the phase variable $\phi$}
	\label{sub:phi}
	
	We begin with a lemma on the sign of $F'$ at the boundary of the interval $[0,1]$,
	which is indispensable for handling the term $F'(\phi)$ in the subsequent proofs.
	
	\begin{lemma}[\textbf{Sign of} $F'$ \textbf{at the boundary of} {$[0,1]$}]
		\label{lem:F_sign}
		With $F'(\phi) = \tfrac12\phi(1-\phi)(1-2\phi)$, we have:
		\begin{itemize}
			\item For $\phi < 0$: $F'(\phi)\cdot\pneg{\phi} \leq 0$, i.e.,
			$F'(\phi)\cdot(-\pneg{\phi}) \geq 0$.
			\item For $\phi > 1$: $F'(\phi)\cdot\ppos{\phi-1} \geq 0$.
		\end{itemize}
	\end{lemma}
	
	\begin{proof}
		For $\phi < 0$, we have $1 - \phi > 0$ and $1 - 2\phi > 0$, so
		$F'(\phi) = \frac12\phi(1-\phi)(1-2\phi) \leq 0$ (since $\phi < 0$),
		and $\pneg{\phi} = -\phi > 0$, whence $F'(\phi)\cdot\pneg{\phi} \leq 0$.\\
		For $\phi > 1$, we have $\phi > 0$, $(\phi-1) > 0$ and $1-2\phi < 0$, hence
		$F'(\phi) = -\frac12\phi(\phi-1)(-(1-2\phi)) = \frac12\phi(\phi-1)(2\phi-1) \geq 0$,
		and $\ppos{\phi-1} = \phi - 1 > 0$, whence $F'(\phi)\cdot\ppos{\phi-1} \geq 0$.
	\end{proof}
	
	\begin{theorem}[\textbf{Confinement of $\phi$ in} {$[0,1]$}]
		\label{thm:phi_bounds}
		Under hypotheses \textup{\textbf{(H1)}}-\textup{\textbf{(H6)}} and assuming that
		$c \in [0,1]$ a.e. in $\QT$ \textup{(Theorem~\ref{thm:c_bounds})}, any solution $\phi$
		of the system~\eqref{eq:sys}-\eqref{eq:CI} satisfies:
		\[
		0 \leq \phi(x,t) \leq 1 \quad\text{a.e. in } \QT.
		\]
	\end{theorem}
	
	\begin{proof}
		\textbf{Lower bound ($\phi \geq 0$).}
		We multiply the third equation of \eqref{eq:sys} by $-\pneg{\phi} \leq 0$ and integrate
		over $\Omega$. The time term gives $\frac{\tau_\phi}{2}\frac{\dd}{\dt}\norm{\pneg{\phi}}_{L^2(\Omega)}^2$.
		The diffusion term gives $\varepsilon^2\norm{\nabla\pneg{\phi}}_{L^2(\Omega)}^2 \geq 0$ after
		integration by parts. The potential term satisfies
		$\int_\Omega F'(\phi)\cdot(-\pneg{\phi})\dx \geq 0$ by Lemma~\ref{lem:F_sign}. For the
		forcing term, on $\{\phi < 0\}$ with $\pneg{\phi} = -\phi > 0$ and $c \geq 0$:
		\[
		\lambda(c-\phi)(-\pneg{\phi})
		= -\lambda c\pneg{\phi} - \lambda[\pneg{\phi}]^2 \leq 0.
		\]
		We therefore obtain:
		\[
		\frac{\tau_\phi}{2}\frac{\dd}{\dt}\norm{\pneg{\phi}(t)}_{L^2(\Omega)}^2
		+ \varepsilon^2\norm{\nabla\pneg{\phi}(t)}_{L^2(\Omega)}^2 \leq 0.
		\]
		The monotone decrease of $t \mapsto \norm{\pneg{\phi}(t)}_{L^2(\Omega)}^2$ and the condition
		$\phi_0 \geq 0$ conclude.
		
		\textbf{Upper bound ($\phi \leq 1$).}
		Set $\psi = \phi - 1$, which satisfies
		$\tau_\phi\,\ptn\psi - \varepsilon^2\Delta\psi + F'(\psi+1) = \lambda(c-\psi-1)$,
		with $\psi(\cdot,0) = \phi_0 - 1 \leq 0$ and homogeneous Neumann conditions. We multiply by
		$\ppos{\psi} \geq 0$. The potential term satisfies
		$\int_\Omega F'(\phi)\ppos{\psi}\dx \geq 0$ by Lemma~\ref{lem:F_sign}. For the
		forcing term, on $\{\psi > 0\} = \{\phi > 1\}$, using $c \leq 1$:
		\[
		\lambda(c-\phi)\ppos{\psi}
		= \lambda(c-1)\ppos{\psi} - \lambda[\ppos{\psi}]^2 \leq 0,
		\]
		since $c - 1 \leq 0$ and $\ppos{\psi} \geq 0$. We obtain:
		\[
		\frac{\tau_\phi}{2}\frac{\dd}{\dt}\norm{\ppos{\psi}(t)}_{L^2(\Omega)}^2
		+ \varepsilon^2\norm{\nabla\ppos{\psi}(t)}_{L^2(\Omega)}^2
		+ \lambda\norm{\ppos{\psi}(t)}_{L^2(\Omega)}^2 \leq 0.
		\]
		The initial condition $\ppos{\psi(\cdot,0)} = 0$ concludes.
	\end{proof}
	
	
	
	
	% ============================================================
	\section{Local asymptotic stability analysis}
	\label{sec:stab}
	% ============================================================
	
	The results of the preceding section ensure the invariance of the domain
	\[
	\mathcal{D} = \{ (\theta,c,\phi) : \theta>0,\; 0\le c\le 1,\; 0\le \phi\le 1 \}
	\]
	on every interval $[0,\mathcal{T}]$ with $\mathcal{T}<T_0$, as well as the existence of a uniform bound
	\[
	\theta(x,t) \ge \underline{\theta}_\mathcal{T} > 0 \quad \text{in } Q_{\mathcal{T}}.
	\]
	These properties allow us to control the nonlinearities of the system, in particular the Arrhenius-type reaction terms.
	
	We are interested in this section in the long-time behaviour of the system in the free regime, that is, in the absence of external forcing ($H_{\mathrm{ext}}=0$). We analyse the stability of a homogeneous stationary state $(\bar{\theta},\bar{c},\bar{\phi})$ by means of an energy method.
	
	
	\subsection{Perturbation system}
	
	We introduce the perturbations
	\[
	u = \theta - \bar\theta, \qquad v = c - \bar c, \qquad w = \phi - \bar\phi.
	\]
	Substituting into \eqref{eq:sys}, the perturbation system reads:
	\begin{equation}
		\label{eq:uvw}
		\left\{
		\begin{aligned}
			\rho c_p \ptn u - \nabla\cdot(k(\phi)\nabla u)
			&= \alpha L_c \ptn v + \alpha L_\phi \ptn w, \\[4pt]
			\ptn v - d_c \Delta v
			&= \mathcal{R}(\theta,c) - \mathcal{R}(\bar\theta,\bar c), \\[4pt]
			\tau_\phi \ptn w - \varepsilon^2\Delta w
			+ (F'(\phi) - F'(\bar\phi))
			&= \lambda(v - w),
		\end{aligned}
		\right.
	\end{equation}
	with homogeneous Neumann boundary conditions, where we have set
	$\mathcal{R}(\theta,c) = -\beta K_d(\theta)\,c + \gamma K_r(\theta)\,(1-c)$.
	
	\subsection{Relative energy functional and coercivity}
	
	We associate to the perturbations $(u, v, w)$ the relative energy functional:
	\begin{equation}
		\label{eq:energy}
		\begin{aligned}
			\mathcal{E}(t) & = \;
			\frac{\rho c_p}{2}\norm{u(t)}_{L^2(\Omega)}^2
			+ \frac12\norm{v(t)}_{L^2(\Omega)}^2
			+ \frac{\tau_\phi}{2}\norm{w(t)}_{L^2(\Omega)}^2
			+ \frac{\varepsilon^2}{2}\norm{\nabla w(t)}_{L^2(\Omega)}^2 \\
			&\quad + \int_\Omega G(w(t))\dx
			- \alpha L_c \int_\Omega u\,v\dx
			- \alpha L_\phi \int_\Omega u\,w\dx,
		\end{aligned}
	\end{equation}
	where $G(z) = F(\bar\phi + z) - F(\bar\phi) - F'(\bar\phi)\,z$ is the residual quadratic part
	of the potential $F$ near $\bar\phi$.
	
	\begin{lemma}[\textbf{Local coercivity of the energy}]
		\label{lem:coercive}
		If $\mathcal T<T_0$, there exist $\delta_0 > 0$ and $C_1, C_2 > 0$ such that
		if $\norm{w}_{L^\infty(\Omega)} \leq \delta_0$ and $\alpha \in [0,\alpha_0)$, where $\alpha_0$ is defined in \eqref{eq:alpha0}, then
		\[
		C_1\bigl(\norm{u}_{L^2(\Omega)}^2 + \norm{v}_{L^2(\Omega)}^2 + \norm{w}_{H^1(\Omega)}^2\bigr)
		\leq \mathcal{E}(t)
		\leq C_2\bigl(\norm{u}_{L^2(\Omega)}^2 + \norm{v}_{L^2(\Omega)}^2 + \norm{w}_{H^1(\Omega)}^2\bigr).
		\]
	\end{lemma}
	
	\begin{proof}
		Assume $\mathcal T<T_0$. Since $F \in C^2$ and $F''(\bar\phi) > 0$, there exist $\delta_0, m_F, M_F > 0$ such that
		$m_F z^2 \leq G(z) \leq M_F z^2$ for $|z| \leq \delta_0$. By integration,
		$m_F\norm{w}_{L^2(\Omega)}^2 \leq \int_\Omega G(w)\dx \leq M_F\norm{w}_{L^2(\Omega)}^2$.
		The cross terms are controlled by Young's inequality:
		\[
		\abs{\alpha L_c \int_\Omega u\,v\dx}
		\leq \frac{\rho c_p}{4}\norm{u}_{L^2(\Omega)}^2 + C\alpha^2\norm{v}_{L^2(\Omega)}^2,
		\quad
		\abs{\alpha L_\phi \int_\Omega u\,w\dx}
		\leq \frac{\rho c_p}{4}\norm{u}_{L^2(\Omega)}^2 + C\alpha^2\norm{w}_{L^2(\Omega)}^2.
		\]
		 For $\alpha < \alpha_0$, the conditions
		$\dfrac{\alpha^2 L_c^2}{\rho c_p} < \dfrac{1}{2}$
		and
		$\dfrac{\alpha^2 L_\phi^2}{\rho c_p} < \dfrac{2m_F + \tau_\phi}{2}$
		are satisfied, and the cross terms are absorbed into the diagonal terms.
	\end{proof}
	
	\subsection{Lipschitz continuity of the reaction term}
	
	\begin{lemma}[\textbf{Lipschitz control of} $\mathcal{R}$]
		\label{lem:lipschitzR}
		Assume that $\mathcal T<T_0$. The function
		$$\mathcal{R}(\theta,c) = -\beta K_d(\theta)\,c + \gamma K_r(\theta)\,(1-c)$$
		is Lipschitz continuous in $(\theta, c)$ on $[\underline\theta_{\mathcal{T}}, \infty) \times [0,1]$: there
		exists $L_R > 0$ such that
		\[
		|\mathcal{R}(\theta,c) - \mathcal{R}(\bar\theta,\bar c)|
		\leq L_R\bigl(|u| + |v|\bigr)
		\quad\text{a.e. in } Q_T.
		\]
	\end{lemma}
	
	\begin{proof}
		Assume $\mathcal T<T_0$. The functions $K_d$ and $K_r$ are of class $C^1$ on $(0,\infty)$, hence Lipschitz continuous
		on every compact subset of $[\underline\theta_{\mathcal{T}}, \infty)$. Since $c, \bar c \in [0,1]$, the conclusion
		follows by a direct computation using the triangle inequality on each term.
	\end{proof}
	
	\subsection{Local asymptotic stability}
	
	\begin{theorem}[\textbf{Local asymptotic stability}]
		\label{thm:stability}
		Under hypotheses \textup{\textbf{(H7)}}-\textup{\textbf{(H10)}}, there exist $\delta > 0$ and $\kappa > 0$ such
		that if
		\[
		\norm{u_0}_{L^2(\Omega)} + \norm{v_0}_{L^2(\Omega)} + \norm{w_0}_{H^1(\Omega)} \leq \delta,
		\]
		then the corresponding solution satisfies, for all $t \in [0, T]$:
		\[
		\mathcal{E}(t) \leq \mathcal{E}(0)\,e^{-\kappa t}.
		\]
		In particular,
		\[
		\norm{u(t)}_{L^2(\Omega)}^2 + \norm{v(t)}_{L^2(\Omega)}^2 + \norm{w(t)}_{H^1(\Omega)}^2
		\leq C\,e^{-\kappa t},
		\]
		and therefore $(\theta, c, \phi) \to (\bar\theta, \bar c, \bar\phi)$ as $t \to +\infty$.
	\end{theorem}
	
	\begin{proof}
		The proof proceeds in seven steps. For simplicity, we adopt the notation $\norm{\cdot} = \norm{\cdot}_{L^2(\Omega)}$.
		
		\medskip
		\textbf{Step 1: Thermal estimate.}
		We test the first equation of \eqref{eq:uvw} by $u$:
		\begin{equation}
			\label{eq3}
			\frac{\rho c_p}{2}\frac{d}{dt}\norm{u}^2 + k_*\norm{\nabla u}^2
			\leq \alpha L_c \int_\Omega \ptn v\,u\dx + \alpha L_\phi \int_\Omega \ptn w\,u\dx,
		\end{equation}
		using $k(\phi) \geq k_* > 0$ by \textbf{(H1)}.
		
		\medskip
		\textbf{Step 2: Chemical estimate.}
		We test the second equation of \eqref{eq:uvw} by $v$:
		\[
		\frac12\frac{d}{dt}\norm{v}^2 + d_c\norm{\nabla v}^2
		= \int_\Omega (\mathcal{R}(\theta,c) - \mathcal{R}(\bar\theta,\bar c))\,v\dx.
		\]
		By Lemma~\ref{lem:lipschitzR} and Young's inequality ($\eta > 0$ to be chosen):
		\begin{equation}
			\label{eq:v-est}
			\frac12\frac{d}{dt}\norm{v}^2 + d_c\norm{\nabla v}^2
			\leq \eta\norm{u}^2 + C_\eta\norm{v}^2.
		\end{equation}
		
		\medskip
		\textbf{Step 3: Phase estimate.}
		We test the third equation of \eqref{eq:uvw} by $w$:
		\[
		\frac{\tau_\phi}{2}\frac{d}{dt}\norm{w}^2 + \varepsilon^2\norm{\nabla w}^2
		+ \int_\Omega (F'(\phi)-F'(\bar\phi))\,w\dx
		= \lambda\int_\Omega v\,w\dx - \lambda\norm{w}^2.
		\]
		Under hypothesis $F''(\bar\phi) > 0$, there exists $m_F > 0$ such that for $w$ sufficiently small,
		$\int_\Omega (F'(\phi)-F'(\bar\phi))\,w\dx \geq m_F\norm{w}^2$. By Young's inequality:
		\begin{equation}
			\label{eq:w-est}
			\frac{\tau_\phi}{2}\frac{d}{dt}\norm{w}^2 + \varepsilon^2\norm{\nabla w}^2
			+ \left(m_F + \frac\lambda2\right)\norm{w}^2
			\leq \frac\lambda2\norm{v}^2.
		\end{equation}
		
		\medskip
		\noindent\textbf{Step 4: Time derivative of the cross terms and cancellation.}
		
		We compute the time derivative of the two cross terms of
		$\mathcal{E}(t)$ defined in \eqref{eq:energy}.
		
		\begin{itemize}
			\item[\textbf{(i)}]
			By Leibniz's rule, we have
			\begin{equation}
				\label{eq6}
				\frac{d}{dt}\!\left(-\alpha L_c\int_\Omega u\,v\dx\right)
				= -\alpha L_c\int_\Omega (\ptn u)\,v\dx
				-\alpha L_c\int_\Omega u\,(\ptn v)\dx.
			\end{equation}
			The second term (s.t.) of \eqref{eq6} exactly cancels the contribution
			$+\alpha L_c\int_\Omega (\ptn v)\,u\dx$ appearing on the right-hand side
			of \eqref{eq3}:
			\[
			\underbrace{+\alpha L_c\int_\Omega (\ptn v)\,u\dx}_{\text{Step 1}}
			+
			\underbrace{\left( -\alpha L_c\int_\Omega u\,(\ptn v)\dx \right) }_{\text{s.t. of (4.6)}} 
			= 0.
			\]
			It remains to handle $-\alpha L_c\int_\Omega (\ptn u)\,v\dx$.
			We extract $\ptn u$ from the first equation of \eqref{eq:uvw}:
			\[
			\rho c_p\,\ptn u
			= \nabla\cdot\bigl(k(\phi)\nabla u\bigr)
			+ \alpha L_c\,\ptn v
			+ \alpha L_\phi\,\ptn w,
			\]
			whence, after dividing by $\rho c_p$:
			\begin{align}
				\label{eq7}
				\begin{split}
					-\alpha L_c\int_\Omega (\ptn u)\,v\dx
					&= -\frac{\alpha L_c}{\rho c_p}
					\int_\Omega \nabla\cdot\bigl(k(\phi)\nabla u\bigr)\,v\dx \\
					&\quad
					-\frac{\alpha^2 L_c^2}{\rho c_p}
					\int_\Omega (\ptn v)\,v\dx
					-\frac{\alpha^2 L_c L_\phi}{\rho c_p}
					\int_\Omega (\ptn w)\,v\dx.
				\end{split}
			\end{align}
			For the first term on the right-hand side of \eqref{eq7}, integration by parts under homogeneous
			Neumann conditions gives
			\[
			-\int_\Omega \nabla\cdot\bigl(k(\phi)\nabla u\bigr)\,v\dx
			= \int_\Omega k(\phi)\,\nabla u\cdot\nabla v\dx,
			\]
			and Young's inequality with $\varepsilon_1 > 0$, using
			$k(\phi) \leq k^*$ by \textbf{(H1)}, yields
			\[
			\frac{\alpha L_c}{\rho c_p}
			\abs{\int_\Omega k(\phi)\,\nabla u\cdot\nabla v\dx}
			\leq
			\frac{\alpha L_c\, k^*}{\rho c_p}\norm{\nabla u}\,\norm{\nabla v}
			\leq
			\frac{\varepsilon_1}{2}\norm{\nabla u}^2
			+ \frac{\alpha^2 L_c^2\,(k^*)^2}{2\varepsilon_1\,\rho^2 c_p^2}\norm{\nabla v}^2.
			\]
			For the second term on the right-hand side of \eqref{eq7}, we have
			\[
			-\frac{\alpha^2 L_c^2}{\rho c_p}
			\int_\Omega (\ptn v)\,v\dx
			= -\frac{\alpha^2 L_c^2}{2\rho c_p}\frac{d}{dt}\norm{v}^2,
			\]
			which is absorbed into the term $\dfrac{1}{2}\dfrac{d}{dt}\norm{v}^2$ from \eqref{eq:v-est} since $\alpha < \alpha_0$ implies
			$\dfrac{\alpha^2 L_c^2}{\rho c_p} < \dfrac{1}{2} < 1$.
			For the third term on the right-hand side of \eqref{eq7}, Young's inequality gives
			\[
			\frac{\alpha^2 L_c L_\phi}{\rho c_p}
			\abs{\int_\Omega (\ptn w)\,v\dx}
			\leq
			\frac{\alpha^2 L_c L_\phi}{2\rho c_p}
			\bigl(\norm{\ptn w}^2 + \norm{v}^2\bigr),
			\]
			terms which will be absorbed into the dissipative contributions of
			\eqref{eq:w-est} and into $\mu_4\norm{w}^2$ from Step~5 for small $\alpha$.
			\item[\textbf{(ii)}]
			By a strictly analogous computation, we find
			\begin{equation}
				\label{eq8}
				\frac{d}{dt}\!\left(-\alpha L_\phi\int_\Omega u\,w\dx\right)
				= -\alpha L_\phi\int_\Omega (\ptn u)\,w\dx
				-\alpha L_\phi\int_\Omega u\,(\ptn w)\dx.
			\end{equation}
			The second term on the right-hand side of \eqref{eq8} exactly cancels the contribution
			$+\alpha L_\phi\int_\Omega (\ptn w)\,u\dx$ from \eqref{eq3} in Step 1:
			\[
			\underbrace{+\alpha L_\phi\int_\Omega (\ptn w)\,u\dx}_{\text{Step 1}}
			+
			\underbrace{\left( -\alpha L_\phi\int_\Omega u\,(\ptn w)\dx \right)}_{\text{s.t. of (4.8)}}
			= 0.
			\]
			For the remaining term $-\alpha L_\phi\int_\Omega (\ptn u)\,w\dx$,
			we again substitute $\rho c_p\,\ptn u$ from \eqref{eq:uvw} to obtain
			\begin{align*}
				-\alpha L_\phi\int_\Omega (\ptn u)\,w\dx
				&= -\frac{\alpha L_\phi}{\rho c_p}
				\int_\Omega \nabla\cdot\bigl(k(\phi)\nabla u\bigr)\,w\dx \\
				&\quad
				-\frac{\alpha^2 L_c L_\phi}{\rho c_p}
				\int_\Omega (\ptn v)\,w\dx
				-\frac{\alpha^2 L_\phi^2}{\rho c_p}
				\int_\Omega (\ptn w)\,w\dx.
			\end{align*}
			After integration by parts and Young's inequality with
			$\varepsilon_2 > 0$, we obtain
			\[
			\frac{\alpha L_\phi}{\rho c_p}
			\abs{\int_\Omega k(\phi)\,\nabla u\cdot\nabla w\dx}
			\leq
			\frac{\varepsilon_2}{2}\norm{\nabla u}^2
			+\frac{\alpha^2 L_\phi^2\,(k^*)^2}{2\varepsilon_2\,\rho^2 c_p^2}
			\norm{\nabla w}^2.
			\]
			The term involving $\ptn w$ gives
			\[
			-\frac{\alpha^2 L_\phi^2}{\rho c_p}
			\int_\Omega (\ptn w)\,w\dx
			= -\frac{\alpha^2 L_\phi^2}{2\rho c_p}\frac{d}{dt}\norm{w}^2,
			\]
			which is absorbed into $\dfrac{\tau_\phi}{2}\dfrac{d}{dt}\norm{w}^2$ from
			\eqref{eq:w-est} since $\alpha < \alpha_0$ implies
			$\dfrac{\alpha^2 L_\phi^2}{\rho c_p} < \tau_\phi$.
			The remaining cross term is handled by Young's inequality:
			\[
			\frac{\alpha^2 L_c L_\phi}{\rho c_p}
			\abs{\int_\Omega (\ptn v)\,w\dx}
			\leq
			\frac{\alpha^2 L_c L_\phi}{2\rho c_p}
			\bigl(\norm{\ptn v}^2 + \norm{w}^2\bigr).
			\]
			
			\item[\textbf{(iii)}]
			Gathering the contributions from Steps~(i) and~(ii), we obtain,
			for all $\varepsilon_1, \varepsilon_2 > 0$ and $\alpha$ sufficiently small, that
			\begin{align*}
				\frac{d}{dt}\!\left(
				-\alpha L_c\int_\Omega uv\dx
				-\alpha L_\phi\int_\Omega uw\dx
				\right)
				&\leq
				(\varepsilon_1+\varepsilon_2)\norm{\nabla u}^2
				+ C_{\varepsilon_1}\alpha^2\norm{\nabla v}^2
				+ C_{\varepsilon_2}\alpha^2\norm{\nabla w}^2 \\
				&\quad + C\alpha^2\bigl(\norm{v}^2 + \norm{w}^2\bigr),
			\end{align*}
			where the constants $C_{\varepsilon_i} > 0$ depend only on
			$L_c, L_\phi, k^*, \rho, c_p$.
			We choose $\varepsilon_1 = \varepsilon_2 = \dfrac{k_*}{4}$, so that
			$\varepsilon_1 + \varepsilon_2 = \dfrac{k_*}{2}$, in order to absorb these terms
			into $k_*\norm{\nabla u}^2$ from Step~1. For $\alpha < \alpha_0$, the conditions
			\[
			\frac{\alpha^2 L_c^2\,(k^*)^2}{2\varepsilon_1\,\rho^2 c_p^2} \leq \frac{d_c}{2}
			\qquad\text{and}\qquad
			\frac{\alpha^2 L_\phi^2\,(k^*)^2}{2\varepsilon_2\,\rho^2 c_p^2} \leq \frac{\varepsilon^2}{2}
			\]
			are satisfied, and the contributions of order $\alpha^2$ are absorbed into
			$d_c\norm{\nabla v}^2$ and $\varepsilon^2\norm{\nabla w}^2$ from Steps~2 and~3. This prepares the global energy inequality of Step~5.
		\end{itemize}
		\textbf{Step 5: Global energy inequality.}
		Summing the estimates from Steps~1 to~4, with $\varepsilon_1 = \varepsilon_2 = \dfrac{k_*}{4}$,
		$\eta \leq \dfrac{\mu_1 C_P}{2}$ and $\alpha < \alpha_0$, we obtain:
		\[
		\frac{d}{dt}\mathcal{E}(t)
		+ \mu_1\norm{\nabla u}^2 + \mu_2\norm{\nabla v}^2 + \mu_3\norm{\nabla w}^2 + \mu_4\norm{w}^2
		\leq C\bigl(\norm{u}^2 + \norm{v}^2\bigr),
		\]
		with constants $\mu_i > 0$.
		
		\medskip
		\textbf{Step 6: Closure by Poincaré's inequality.}
		Since $\int_\Omega u\dx = 0$ for all time (conservation of the thermal mean under homogeneous Neumann conditions and $H_{\mathrm{ext}} = 0$), Poincaré's inequality
		then gives $\norm{u}^2 \leq C_P\norm{\nabla u}^2$. For $v$, the reaction term controls
		$\norm{v}^2$ via the stability of $\bar c$. Combining with the coercivity of $\mathcal{E}$:
		\[
		\frac{d}{dt}\mathcal{E}(t) + \kappa\,\mathcal{E}(t) \leq 0
		\]
		for some constant $\kappa > 0$.
		
		\medskip
		\textbf{Step 7: Conclusion by Gronwall.}
		The Gronwall lemma applied to $\psi(t) = \mathcal{E}(t)$ gives
		$\mathcal{E}(t) \leq \mathcal{E}(0)\,e^{-\kappa t}$.
		The equivalence of $\mathcal{E}(t)$ with the norm $\norm{u}_{L^2(\Omega)}^2 + \norm{v}_{L^2(\Omega)}^2 + \norm{w}_{H^1(\Omega)}^2$ completes the proof.
	\end{proof}
	
	\begin{remark}
		The decay rate $\kappa$ depends explicitly on $k_*$, $d_c$, $\varepsilon^2$,
		$\lambda$, $m_F$, and the Poincaré constant $C_P$ of the domain $\Omega$. It is therefore
		amenable to quantitative estimation once the model parameters are fixed. For the
		potential \eqref{eq:F}, one has $F''(\phi) = 1 - 6\phi + 6\phi^2$, so that the condition
		$F''(\bar\phi) > 0$ is equivalent to $\bar\phi \notin \left(\frac{3-\sqrt{3}}{6},\frac{3+\sqrt{3}}{6}\right)$,
		which excludes the neighbourhood of the local maximum of $F$ at $\phi = \tfrac{1}{2}\cdotp$
	\end{remark}
	
	\section{Conclusion}
	In this work, we have studied a strongly coupled thermo-reaction-phase system modelling the thermal degradation of heat-sensitive materials. Despite the presence of singular Arrhenius-type nonlinearities and couplings involving time derivatives, we have established a maximum principle ensuring the positivity of the temperature as well as the invariance of the physically admissible domain. We have then proved the local asymptotic stability of a homogeneous stationary state in the free regime, by means of a suitably chosen relative energy functional, leading to exponential decay of perturbations. Finally, the proposed system provides a mathematical framework for describing thermally induced degradation processes. More generally, it belongs to a broader class of coupled diffusion--reaction--phase systems, where the phase variable may represent different macroscopic quantities depending on the context. In this work, it is interpreted as a measure of visibility, allowing one to relate this phenomenon to physical parameters and environmental conditions. The present analysis establishes the theoretical foundations for subsequent quantitative studies, such as sensitivity analysis and parameter optimization, which will be addressed in future work.
	

	
	\section*{References}
	
\begin{thebibliography}{00}
	
	\bibitem{AllenCahn1979}
	S.M. Allen and J.W. Cahn,
	\textit{A microscopic theory for antiphase boundary motion and its application
		to antiphase domain coarsening},
	Acta Metallurgica, vol.~27, no.~6, pp.~1085--1095, 1979.
	\newblock \url{https://doi.org/10.1016/0001-6160(79)90196-2}
	
	\bibitem{BebernesEberly1989}
	J. Bebernes and D. Eberly,
	\textit{Mathematical Problems from Combustion Theory},
	Applied Mathematical Sciences, vol.~83, Springer, New York, 1989.
	\newblock \url{https://doi.org/10.1007/978-1-4612-4546-9}
	
	\bibitem{Caginalp1986}
	G. Caginalp,
	\textit{An analysis of a phase field model of a free boundary},
	Archive for Rational Mechanics and Analysis, vol.~92, no.~3,
	pp.~205--245, 1986.
	\newblock \url{https://doi.org/10.1007/BF00254827}
	
	\bibitem{CahnHilliard1958}
	J.W. Cahn and J.E. Hilliard,
	\textit{Free energy of a nonuniform system. I. Interfacial free energy},
	The Journal of Chemical Physics, vol.~28, no.~2, pp.~258--267, 1958.
	\newblock \url{https://doi.org/10.1063/1.1744102}
	
	\bibitem{ColliSprekels1995}
	P. Colli and J. Sprekels,
	\textit{On a Penrose--Fife model with zero interfacial energy leading to a
		phase-field system of relaxed Stefan type},
	Annali di Matematica Pura ed Applicata, vol.~169, no.~1,
	pp.~269--289, 1995.
	\newblock \url{https://doi.org/10.1007/BF01759357}
	
	\bibitem{ElliottZheng1986}
	C.M. Elliott and S. Zheng,
	\textit{On the Cahn--Hilliard equation},
	Archive for Rational Mechanics and Analysis, vol.~96, no.~4,
	pp.~339--357, 1986.
	\newblock \url{https://doi.org/10.1007/BF00251803}
	
	\bibitem{EvansBook2010}
	L.C. Evans,
	\textit{Partial Differential Equations},
	2\textsuperscript{nd}~ed., Graduate Studies in Mathematics, vol.~19,
	American Mathematical Society, Providence, 2010.
	
	\bibitem{EvansEtal1992}
	L.C. Evans, H.M. Soner and P.E. Souganidis,
	\textit{Phase transitions and generalized motion by mean curvature},
	Communications on Pure and Applied Mathematics, vol.~45, no.~9,
	pp.~1097--1123, 1992.
	\newblock \url{https://doi.org/10.1002/cpa.3160450903}
	
	\bibitem{KenmochiNiezgodka1994}
	N. Kenmochi and M. Niezg\'odka,
	\textit{Systems of nonlinear parabolic equations for phase change problems},
	Advances in Mathematical Sciences and Applications, vol.~3,
	pp.~89--117, 1993--1994.
	
	\bibitem{MacLarenWhite2005}
	D.C. MacLaren and M.A. White,
	\textit{Design rules for reversible thermochromic mixtures},
	Journal of Materials Science, vol.~40, no.~3, pp.~669--676, 2005.
	\newblock \url{https://doi.org/10.1007/s10853-005-6109-x}
	
	\bibitem{MakkiMiranvilleSaoud2020}
	A. Makki, A. Miranville and W. Saoud,
	\textit{On a Cahn--Hilliard/Allen--Cahn system coupled with a type~{III}
		heat equation and singular potentials},
	Nonlinear Analysis: Theory, Methods and Applications, vol.~196,
	article~111804, 2020.
	\newblock \url{https://doi.org/10.1016/j.na.2020.111804}
	
	\bibitem{MiranvilleBook2019}
	A. Miranville,
	\textit{The Cahn--Hilliard Equation: Recent Advances and Applications},
	CBMS-NSF Regional Conference Series in Applied Mathematics, vol.~95,
	SIAM, Philadelphia, 2019.
	\newblock \url{https://doi.org/10.1137/1.9781611975925}
	
	\bibitem{MiranvilleQuintanillaSaoud2020}
	A. Miranville, R. Quintanilla and W. Saoud,
	\textit{Asymptotic behavior of a Cahn--Hilliard/Allen--Cahn system
		with temperature},
	Communications on Pure and Applied Analysis, vol.~19, no.~4,
	pp.~2257--2288, 2020.
	\newblock \url{https://doi.org/10.3934/cpaa.2020099}
	
	\bibitem{MiranvilleZelik2004}
	A. Miranville and S. Zelik,
	\textit{Robust exponential attractors for Cahn--Hilliard type equations
		with singular potentials},
	Mathematical Methods in the Applied Sciences, vol.~27, no.~5,
	pp.~545--582, 2004.
	\newblock \url{https://doi.org/10.1002/mma.464}
	
	\bibitem{Ntsokongo2023}
	A.J. Ntsokongo,
	\textit{Asymptotic behavior of an Allen--Cahn type equation with temperature},
	Discrete and Continuous Dynamical Systems -- S, vol.~16, no.~9,
	pp.~2452--2466, 2023.
	\newblock \url{https://doi.org/10.3934/dcdss.2023017}
	
	\bibitem{PenroseFife1990}
	O. Penrose and P.C. Fife,
	\textit{Thermodynamically consistent models of phase-field type for the
		kinetics of phase transitions},
	Physica D: Nonlinear Phenomena, vol.~43, no.~1, pp.~44--62, 1990.
	\newblock \url{https://doi.org/10.1016/0167-2789(90)90015-H}
	
	\bibitem{PierreM2010}
	M. Pierre,
	\textit{Global existence in reaction-diffusion systems with control of mass:
		a survey},
	Milan Journal of Mathematics, vol.~78, no.~2, pp.~417--455, 2010.
	\newblock \url{https://doi.org/10.1007/s00032-010-0133-4}
	
	\bibitem{SeebothEtal2014}
	A. Seeboth, D. L\"otzsch, R. Ruhmann and O. Muehling,
	\textit{Thermochromic Polymers --- Function by Design},
	Chemical Reviews, vol.~114, no.~5, pp.~3037--3068, 2014.
	\newblock \url{https://doi.org/10.1021/cr400462e}
	
	\bibitem{StuartEtal2010}
	M.A.C. Stuart, W.T.S. Huck, J. Genzer et al.,
	\textit{Emerging applications of stimuli-responsive polymer materials},
	Nature Materials, vol.~9, no.~2, pp.~101--113, 2010.
	\newblock \url{https://doi.org/10.1038/nmat2614}
	
	\bibitem{TemamBook1997}
	R. Temam,
	\textit{Infinite-Dimensional Dynamical Systems in Mechanics and Physics},
	2\textsuperscript{nd}~ed., Applied Mathematical Sciences, vol.~68,
	Springer, New York, 1997.
	\newblock \url{https://doi.org/10.1007/978-1-4612-0645-3}
	
	\bibitem{Ziemer}
	W.P. Ziemer,
	\textit{Weakly Differentiable Functions},
	Springer-Verlag, New York, 1989.
	
\end{thebibliography}

\end{document}

\endinput
%%
%% End of file `elsarticle-template-num-names.tex'.
